% !PRINTED 738
\weberpage{738}

Zur besseren Uebersicht setzen wir die Gleichungen (13)
noch einmal ohne den Index $k$, der die conjugirten Werthe von
einander unterscheidet, hierher:
\weq{14}{
\begin{aligned}
2\, l_i & = m_{1,i}\, l'_1 + m_{2,i}\, l'_2 + \cdots + m_{\nu'-1,i}\, l'_{\nu'-1} \\
& \quad + M_{0,i}\, L_0 + M_{1,i}\, L_1 + M_{2,i}\, L_2 + \cdots + M_{\nu'-1,i}\, L_{\nu'-1}.
\end{aligned}
}

Bilden wir nun nach (13) den Regulator $E$ und bezeichnen
mit $M$ den absoluten Werth der Determinante der $\nu-1$ Reihen
der Zahlen $m_{k,i}$, $M_{k,i}$:
\[
M = \pm \Sigma \pm m_{1,1} \ldots m_{\nu'-1,\nu'-1}\, M_{0,0}\, M_{1,1} \ldots M_{\nu'-1,\nu'-1},
\]
so folgt:
\weq{15}{2^{\nu-1} E = M R,}
und daraus mit Rücksicht auf (12):
\weq{16}{E = 2^{-\nu'} M E' T_0.}

Es lässt sich nun weiter nachweisen, dass $M$ eine Potenz
von 2 und durch $2^{\nu'}$ theilbar ist.

Da nämlich das System (1) als Fundamentalsystem in $H_m$
vorausgesetzt war, so giebt es \gesperrt{ganze rationale Zahlen}
$n_{k,i}$, $N_{k,i}$, die den Gleichungen genügen:
\weq{17}{
\begin{aligned}
l'_i & = n_{1,i}\, l_1 + n_{2,i}\, l_2 + \cdots + n_{\nu-1,i}\, l_{\nu-1}, \quad i = 1, 2, \ldots, \nu'-1 \\
L_i & = N_{1,i}\, l_1 + N_{2,i}\, l_2 + \cdots + N_{\nu-1,i}\, l_{\nu-1}, \quad i = 0, 1, 2, \ldots, \nu'-1.
\end{aligned}
}

Bilden wir daraus die Determinante $R$ und bezeichnen mit $N$
den absoluten Werth der Determinante der $n_{k,i}$, $N_{k,i}$, so folgt:
\[
R = N E,
\]
und daraus nach (15):
\[
M N = 2^{\nu-1};
\]
daraus folgt, dass jede der beiden ganzen Zahlen $M$, $N$ eine
Potenz von 2 sein muss.

Um zu entscheiden, durch welche Potenz von 2 die Determinante
 $M$ theilbar ist, bestimmen wir zunächst ein System
ganzer rationaler Zahlen $a_i$ ohne gemeinschaftlichen Theiler, das
den linearen Gleichungen
\weq{18}{\sum\limits_{1,\, \nu-1}^{i} a_i\, m_{k,i} = 0 \quad (k = 1, 2, \ldots, \nu'-1)}
genügt. Setzen wir dann
\[
\sum\limits_{1,\, \nu-1}^{i} a_i\, M_{k,i} = \xi_k,
\]
%CONT>>
